%-------------------------------------------------------------------------------
%	SECTION TITLE
%-------------------------------------------------------------------------------
\cvsection{Projects}


%-------------------------------------------------------------------------------
%	CONTENT
%-------------------------------------------------------------------------------
\begin{cventries}

%---------------------------------------------------------
  \cventry
    {Course Project} % Role
    {Design of a stimulation based rehabilitative device for gait correction in patients suffering from neurological disorders} % Title
    {University at Buffalo} % Location
    {Aug. 2016 - Dec. 2016} % Date(s)
    {
      \begin{cvitems} % Description(s)
        \item {Designed a rehabilitative device that is able to predict the concurrent phase in the gait cycle in patients suffering from a neurological disorder that causes irregular gait. Using these predictions, the system stimulates the designated muscles to correct deviations in gait cycle arising from their neurological ailment.}
        \item {The system and all inclusive components were presented with a \ul{systems engineering plan (SEP)}.}
      \end{cvitems}
    }
    
%---------------------------------------------------------
  \cventry
    {Course Project} % Role
    {Fabrication of asymmetric lattice-shaped microchannel structures for continuous size-dependent cell sorting} % Title
    {University at Buffalo} % Location
    {Jan. 2016 - May 2016} % Date(s)
    {
      \begin{cvitems} % Description(s)
        \item {By creating a specialized mask via \ul{photolithography} to create a master on a silicone wafer, a \ul{bioMEMs} device was produced from poly(dimethylsiloxane) that contained a mucrofluidic system for cell sorting. Key efficiency parameters were analyzed to optimize the micro-device even further.}
        \item {Learned about \ul{wafer fabrication} and the process of creating bioMEMs to improve healthcare.}
      \end{cvitems}
    }

%---------------------------------------------------------
  \cventry
    {Senior Design Project} % Role
    {Design of an animal positioning system for simultaneous Positron Emission Tomography (PET) and Magnetic Resonance Imaging (MRI) at ultra-high magnetic field strength (Bruker 7T)} % Title
    {University at Buffalo} % Location
    {Aug. 2015 - May 2016} % Date(s)
    {
      \begin{cvitems} % Description(s)
        \item {The combination of PET and MRI modalities allow simultaneous acquisition of high resolution anatomical data using MRI and quantitative physiological data using PET.}
        \item {A supporting bed system was designed which is compatible with both MRI and PET. The apparatus allowed the preparation the animal for the experiment outside of the magnet which then slid into the magnet and into the center of the detector. The system was non-magnetic and non-conducting due to the presence of a very strong magnetic field, supplied vital conditions to the animal, and contained potential radiotracer spills.}
        \item {This system was designed to speed up \ul{simultaneous MRI/PET} image acquisition in mouse studies, to improve imaging throughput.}
      \end{cvitems}
    }

%---------------------------------------------------------
  \cventry
    {Course Project} % Role
    {Characterization of the variation of saccade speeds across the three dimensional field of view in humans} % Title
    {University at Buffalo} % Location
    {Aug. 2014 - Dec. 2014} % Date(s)
    {
      \begin{cvitems} % Description(s)
        \item {Using the BioPac software and hardware suite, and a complex system to track eye position, saccade speeds at different angular positions were acquired to form a three dimensional plot of saccade variations of a subject’s eye when looking at different points in a spatial plane. The plots were then normalized and an equation was derived to predict the subject's saccade speeds at various angles.}
        \item {Inculcated the importance of normalizing data and learned various  statistical approaches to present large data.}
      \end{cvitems}
    }

%---------------------------------------------------------
\end{cventries}
